% Данный файл распространяетсяя под лицензией CC BY 4.0. Текст лицензии размещён на https://creativecommons.org/licenses/by/4.0/
% (c) Кафедра системного программирования, 2025

\documentclass[a4paper]{article}

\usepackage[a4paper, top=8mm, bottom=8mm, left=8mm, right=8mm]{geometry}

\usepackage{polyglossia}
\setdefaultlanguage[babelshorthands=true]{russian}
\setotherlanguage{english}

\usepackage{fontspec}
\setmainfont{FreeSerif}
\newfontfamily{\russianfonttt}[Scale=0.7]{DejaVuSansMono}

\usepackage[tiny, compact]{titlesec}

\usepackage{titling}
\setlength{\droptitle}{-1cm}
\pretitle{\begin{center}\begin{bfseries}\Large}
\posttitle{\par\end{bfseries}\end{center}}
\preauthor{\begin{center}\normalsize}
\postauthor{\par\end{center}\vspace{-1.8cm}}

\usepackage{hyperref}
\usepackage{bookmark}
\usepackage{csquotes}

\title{Разработка самовоспроизводящегося компилятора}

\author{Балышев Артем Максимович}

\date{}

\begin{document}

\maketitle

\begin{flushright}
    Группа: \emph{\enquote{24.Б72-мм}}

    Кафедра: \emph{кафедра системного программирования}

    Научный руководитель: \emph{Косарев Дмитрий Сергеевич}

    Номер семестра практики: \emph{4}
\end{flushright}


\section{Постановка задачи}

Работа выполняется в рамках проекта «rukaml» по разработке компилятора функционального языка. \textbf{Целью работы является} обеспечение возможности сборки языковых инструментов с помощью \texttt{\large rukaml}.

Конечными пользователями продукта являются другие разработчики, в частности компиляторные инженеры.

Значимость результатов работы заключается в демонстрации возможности развития проекта в качестве полноценного инструмента сборки.

\section{Описание предлагаемого решения}

Предлагается доработать \texttt{\large rukaml} до состояния, в котором он способен собрать простой компилятор подмножества языка C, написанный вручную.


Язык C выбран в силу относительной простоты написания компиляторов для него. Написание нового компилятора C вместо поиска существующего решения и попыток обеспечить его сборку обосновывается тем, что \texttt{\large rukaml} поддерживает лишь сильно ограниченное подмножество языка OCaml.


Реализация обоих компиляторов осуществлялась на языке OCaml, реализация стандартной библиотеки \texttt{\large rukaml} — на языке C.


В \texttt{\large rukaml} кодогенерация реализовывалась для архитектуры AMD64, в компиляторе C — для архитектуры RISC-V 64.


\section{Тестирование}


\begin{itemize}
    \item Написание модульных тестов для отдельных компонентов \texttt{\large rukaml}
    \item Написание системных тестов для \texttt{\large rukaml}, проверяющих работоспособность сразу всех этапов компиляции
    \item Сравнение поведения компилируемых программ при компиляции с помощью \texttt{\large ocamlopt} и с помощью \texttt{\large rukaml}
    \item Проведение ручного тестирования компилятора C, собранного как с помощью \texttt{\large rukaml}, так и с помощью \texttt{\large ocamlopt}
\end{itemize}


Тестовое покрытие, измеренное с помощью утилиты \texttt{\large bisect\_ppx}, составило 79~\% строк.


\section{Заключение}

В ходе выполнения работы получены следующие результаты:


\begin{itemize}
    \item Расширение стандартной библиотеки \texttt{\large rukaml} примитивами, необходимыми для сборки консольных утилит:
          \begin{itemize}
              \item Форматированный вывод
              \item Обработка аргументов командной строки
              \item Работа с файловой системой
          \end{itemize}
    \item Реализация преобразования имён (англ. name mangling) и разрешение псевдонимов для примитивов стандартной библиотеки, необходимые для обеспечения непротиворечивости представления компилируемой программы на этапе кодогенерации
    \item Написание простого компилятора, транслирующего подмножество языка C в ассемблерный код RISC-V 64, и его сборка с помощью \texttt{\large rukaml}
    \item Выступление с докладом на секционном заседании конференции СПбПУ «Современные технологии в теории и практике программирования»
          \begin{itemize}
              \item Тезисы доклада включены в сборник материалов конференции, проиндексированный в РИНЦ
          \end{itemize}
\end{itemize}

Ссылки на пулл-реквесты:
\begin{itemize}
    \item https://github.com/Kakadu/rukaml/pull/26
    \item https://github.com/Kakadu/rukaml/pull/27
    \item https://github.com/Kakadu/rukaml/pull/28
\end{itemize}

Техническая документация: в описании к пулл-реквестам прилагаются описания изменений.

Решение находится на этапе кодревью и рефакторинга.

\end{document}
